\begin{prob}
    % Write the problem here.
    Suppose
    \begin{itemize}
        \item algorithm A takes $n$ milliseconds to run on a problem of size $n$,
        \item algorithm B takes $n^2$ milliseconds to run on a problem of size $n$,
        \item algorithm C takes $2^n$ milliseconds to run on a problem of size $n$,
        \item algorithm D takes $\log_2(n)$ milliseconds to run on a problem of size $n$.
    \end{itemize}
    What is the largest problem size each algorithm can solve in 1 second,
    10 seconds, and 1 minute? That is, fill in the following table:

    \begin{center}
        \begin{tabular}{rccc}
            \toprule
              & 1 sec. & 10 sec. & 1 min \\ \midrule
            A & ?      & ?       & ?     \\
            B & ?      & ?       & ?     \\
            C & ?      & ?       & ?     \\
            D & ?      & ?       & ?     \\ \bottomrule
        \end{tabular}
    \end{center}


    \begin{soln}
        First, we convert the times into milliseconds. One second is 1,000 milliseconds;
        ten seconds is 10,000 milliseconds, and one minute is 60,000 milliseconds.

        Algorithm A can solve a problem of size $n$ in $n$ milliseconds. Hence
        It can solve a problem of size $1,000$ in $1,000$ milliseconds;
        a problem of size $10,000$ in $10,000$ milliseconds, and a problem
        of size $60,000$ in $60,000$ milliseconds.

        Algorithm B can solve a problem of size $n$ in $n^2$ milliseconds. Equivalently, it can solve a problem of size $\sqrt{n}$ in $n$ milliseconds. Therefore, it can solve a problem of size $\sqrt{1,000}$ in 1,000 milliseconds. But $\sqrt{1,000}$ isn't an integer, and problem sizes have to be integers. Hence we take the largest integer
        smaller than $\sqrt{1,000}$ (also known as the \emph{floor}), which
        turns out to be 31. The algorithm can solve a problem of size
        $\sqrt{10,000} = 100$ in 10,000 milliseconds, and a problem of size
        $\lfloor \sqrt{60,000} \rfloor = 244$ in 60,000 milliseconds.

        Algorithm C can solve a problem of size $n$ in $2^n$ milliseconds. Equivalently, it can solve a problem of size $\log_2(n)$ in $n$ milliseconds. Therefore, it can solve a problem of size $\lfloor \log_2{(1,000)} \rfloor = 9$ in 1,000 milliseconds; a problem of size $\lfloor \log_2{(10,000)} \rfloor = 13$
        in 10,000 milliseconds; and a problem of size
        $\lfloor \log_2{(60,000)} \rfloor = 15 $ in 60,000 milliseconds.

        Algorithm D can solve a problem of size $n$ in $\log_2(n)$ milliseconds. Equivalently, it can solve a problem of size $2^n$ in $n$ milliseconds. Therefore, it can solve a problem size of $2^{1,000}$ in 1,000 milliseconds; a problem of size $2^{10,000}$; in 10,000 milliseconds; a problem of size $2^{60,000}$ in 60,000 milliseconds.

        Therefore the filled-in table is:

        \begin{center}
            \begin{tabular}{rccc}
                \toprule
                  & 1 sec.      & 10 sec.      & 1 min        \\ \midrule
                A & 1,000       & 10,000       & 60,000       \\
                B & 31          & 100          & 244          \\
                C & 9           & 13           & 15           \\
                D & $2^{1,000}$ & $2^{10,000}$ & $2^{60,000}$ \\ \bottomrule
            \end{tabular}
        \end{center}


    \end{soln}

\end{prob}
